\subsection{C-DAS/SFERA\@}
Under intervjuer med lokförare framkom det att enkel, tydlig och tillförlitlig kommunikation samt kontext som både är tydlig \emph{och} relevant är central för deras vilja att följa hastighetsanvisningar från C-DAS-OB och annan kommunikation från TMS\@.
Det är därför viktigt att SFERA-protokollet fortsätter att utvecklas och att lokförarens perspektiv är med som en röd tråd i utvecklingsarbetet för att tillgodose de behov som finns.
Att utbildning ges för att säkerställa att alla lokförare förstår och kan utnyttja systemets fulla potential.
Bara under intervjuerna märktes en tydlig skillnad i respondenternas inställning till SFERA och C-DAS i takt med att de fick en större förståelse.
\paragraph{}
\hspace{-13pt}
Trots de positiva resultaten är det viktigt att ifrågasätta och granska vissa aspekter av implementeringen.
Exempelvis framkom det under valideringssessionen att vissa meddelanden var onödiga och gängse kommunikation med GSM-R-telefonsamtal hade varit \emph{enklare} än vissa av de delar som presenteras i analysdelens avsnitt\ref{subsec:sfera-utokat}.
Ett återkommande tema under intervjuerna var lokförarnas oro för att det nya systemet skulle kunna öka deras arbetsbörda, särskilt under stressiga förhållanden.
Detta belyser behovet av att noggrant övervaka och utvärdera arbetsmiljöpåverkan under systemets införande, i synnerhet när kommunikationen byggs ut.
Vidare forskning bör undersöka långsiktiga effekter av C-DAS och SFERA på lokförares arbetsmiljö och stressnivåer samt exakt hur kommunikationen går till när alla system är i reguljär drift.
\paragraph{}
\hspace{-14pt}
Trafikverket borde ta fram utbildningsmaterial om C-DAS och SFERA som behandlar dess syfte, bakgrund, ursprung, omfattning, vilka problem de ämnar lösa och hur tekniken fungerar.
Materialet bör vara anpassat för lokförarnas tekniska kompetens och förankras väl i deras vardagliga arbete för att skapa ett så gott intresse som möjligt.
Utbildningsmaterial och eventuella verktyg kan sedan ges ut till järnvägsföretagen som ansvarar för att utbilda sin personal.
Om inte detta görs och ansvaret istället läggs på järnvägsföretagen riskerar utbildning på området bli spretig och kvaliteten kan skilja sig från företag till företag.
Att Trafikverket tar fram en C-DAS-applikation som fungerar enligt SFERA-protokollet är också en idé som framförallt skulle gynna de företag som kör utan digitala förarverktyg idag.
\subsubsection{Begränsningar}
Det kan finnas begränsningar i studien.
Analys har baserats på underlag från intervjuer och om urvalet inte varit optimalt kan fel slutsatser dragits.
Exempelvis fanns bara en (1) godstågsförare av fem respondenter.
Större representation och insikter från annat än persontågsförare vore önskvärt för att uppnå en jämnare fördelning mellan de två.
Det kan även vara så att de förare som intervjuats har en mer positiv bild av digitala verktyg än lokförarkåren i stort.
Detta beslut togs dock för att få ett större och säkrare datainsamlingsunderlag med djupare samtal och tankegångarna då respondenterna var kända sedan tidigare av författaren.
Insikter från andra branschexperter såsom utvecklare eller praktiker som har erfarenhet av C-DAS eller SFERA skulle kunnat öka förtroendet för studien, även om målet var att studera lokförarens perspektiv i dessa system.
Vidare är SFERA som tidigare nämnt relativt nytt och i viss mån fortfarande under utveckling, även om ett färdigt protokoll funnits sedan 2022.
Detta kan betyda att alla dess funktioner inte är fullständigt integrerade eller presenterade och därför kan det vara svårt för respondenterna att fullt ut förstå tekniken.
Slutligen finns tidsaspekten.
Studien har fokuserat på korttidsresultat, och för att undersöka teknikens långsiktiga påverkan på ett mer övergripande plan krävs ytterligare forskning på effekter, arbetsmiljö och andra faktorer som tekniken studien fokuserar på innebär.
\paragraph{}
\hspace{-15pt}
Även om C-DAS och SFERA inte är tvingande för varken infrastrukturförvaltare eller järnvägsföretag att implementera kommer det sannolikt att få väldigt stor inverkan i järnvägsbranschen och hur trafiken bedrivs.
Om dessa digitala system skulle bli obligatoriska kommer kraven på högkvalitativa och robusta system att bli större vilken skulle kräva ytterligare forskning inom ämnet och hur det på bästa vis tillämpas.